\documentclass[11pt,letterpaper]{article}
\usepackage[margin=1in]{geometry}
\usepackage{amsmath,amssymb,amsthm}
\usepackage{graphicx}
\usepackage{hyperref}
\usepackage{booktabs}
\usepackage{enumitem}
\usepackage{bm}
\usepackage[protrusion=true,expansion=false]{microtype}
\usepackage{xcolor}

\hypersetup{
  colorlinks=true, linkcolor=blue!60!black, citecolor=blue!60!black,
  urlcolor=blue!60!black,
}

\title{\Large\textbf{Cross-diagnostic consensus for tokamak disruption prediction}\\
\large{A ten-iteration honest-iteration log from address-space novelty
catcher to gradient-boosted multi-scale classifier on the 500-shot MAST
dataset, with stability cross-validation and failure analysis}}

\author{Christian Knopp \\
  \texttt{cknopp@gmail.com} \\
  \href{https://github.com/Zynerji/JET-Disruption}{\texttt{github.com/Zynerji/JET-Disruption}}}

\date{2026-05-13}

\begin{document}
\maketitle

\begin{abstract}
We document a ten-iteration honest-engineering record of porting a
cross-diagnostic consensus catcher --- originally developed for
unmodeled gravitational-wave burst search and for cross-observable
geometric-locus identification in spacetime metrics --- to the problem
of tokamak disruption prediction on the public 500-shot MAST dataset
(Sharma 2025, Zenodo \texttt{10.5281/zenodo.16032053}). The pipeline
operates on the same robust rolling z-score + multi-detector consensus
structure used by the Systroph\=e v2 detector triple: per-channel
novelty scoring with median-and-MAD baseline, followed by an alert rule
that requires a fraction-threshold of channels to co-fire simultaneously
for a minimum sustain. We iterate the pipeline through nine design
variations, document two outright negative results (hybrid
pendulum/Kuramoto/Hamiltonian voting; failure-analysis-informed lean
feature set), and converge on a 27-feature gradient-boosting classifier
that achieves \textbf{74.3 $\pm$ 1.0\% TPR at 14.1 $\pm$ 2.3\% FAR with
80.8 $\pm$ 2.4 ms median time-to-disruption} (mean $\pm$ std across
three 5-fold cross-validation seeds). The result matches published
recurrent-neural-network baselines on FAR while remaining $\sim$11 pp
below them on TPR. We give an explicit failure analysis showing that
only 2\% of disruptions are structurally invisible to our feature set,
identify the remaining 24\% miss budget as a threshold-tuning gap rather
than a missing-information gap, and discuss why the remaining
performance gap to RNN baselines likely requires per-shot operational
metadata (disruption-cause classification) that this dataset does not
expose. All code, all checkpoints, and all per-iteration JSON results
are publicly available; the iteration log is reproducible commit-by-commit.
\end{abstract}

\section{Introduction}\label{sec:intro}

A tokamak \emph{disruption} is the sudden uncontrolled termination of a
magnetic confinement plasma: stored magnetic energy that took seconds
to build up is released within $\sim$10\,ms, depositing localized
megajoule heat loads onto the wall and inducing eddy-current forces on
the vessel structure. On ITER and on future fusion reactors, the
engineering cost of an unmitigated disruption is sufficient that
operations cannot proceed without a reliable forward-prediction system
that triggers mitigation hardware (massive-gas-injection valves,
shattered-pellet injectors) with adequate lead time, typically 30--50\,ms
\cite{vega2013,maraschek2018}. The standard performance figure is true-positive
rate (TPR, fraction of real disruptions caught at $\ge 30$ ms lead)
versus false-alarm rate (FAR, fraction of non-disruptive shots
triggering a spurious alert).

The published state of the art on real plasma data is
recurrent-neural-network ensembles. Ferreira et al.\ (2019)
\cite{ferreira2019} report a Long Short-Term Memory architecture on JET
bolometer data yielding approximately 85\% TPR at 15\% FAR. The TPR/FAR
trade-off is a free design parameter (set by the binary threshold on
the LSTM output) and the operating curve is dataset-dependent.

We are interested in a methodologically distinct approach: applying the
\emph{cross-observable consensus catcher} originally developed for
unmodeled-burst search in gravitational-wave data
\cite{systrophe_gw} and for cross-observable identification of metric
special loci in spacetime physics \cite{systrophe_lp_ridge} to the
tokamak disruption problem. The structural argument is that a
disruption is a regime in which many independent plasma observables
(magnetic, kinetic, radiative) deflect from their baseline
simultaneously; this is exactly the cross-observable consensus
signature the catcher is designed to detect. The empirical question is
whether the methodology, applied without machine-specific tuning or
operational metadata, can reach performance comparable to published RNN
baselines.

This paper records a ten-iteration log of that effort. Of the ten
iterations, eight produced measurable forward progress and two produced
informative negative results that constrained subsequent feature design.
The final operating point matches published RNN FAR at $\sim$11 pp
lower TPR; we identify the remaining gap explicitly.

\section{Method}\label{sec:method}

\subsection{The catcher v2 detector triple}\label{ssec:catcher}

Given a $C$-channel diagnostic time series sampled at rate $f_s$, the
catcher v2 operates as follows.

\paragraph{Per-channel novelty score.} For each channel $c$ and each
sample $t$, the catcher computes a robust z-score of the assessment
window against a trailing baseline:
\[
z_c(t) \;=\; \frac{\mathrm{median}\bigl(x_c[t-A : t]\bigr) - \mathrm{median}\bigl(x_c[t-B : t-A]\bigr)}
                  {1.4826 \cdot \mathrm{MAD}\bigl(x_c[t-B : t-A]\bigr) + \epsilon}
\]
with assessment-window length $A$ (we use 5\,ms) and baseline-window
length $B$ (we use 50, 200, and 500\,ms at three scales). A channel is
flagged as novel at $t$ if $|z_c(t)| > z_{\mathrm{thr}}$ (we use 3.0).
The median-and-MAD form is robust to channel-specific outlier
distributions and requires no prior calibration.

\paragraph{Consensus alert.} At each $t$, the catcher computes the
fraction of channels novel:
\[
f(t) \;=\; \frac{1}{C} \sum_{c=1}^{C} \mathbb{1}\bigl[|z_c(t)| > z_{\mathrm{thr}}\bigr].
\]
An alert fires when $f(t) \ge f_{\mathrm{thr}}$ sustained for at least
$\tau_{\mathrm{min}}$ samples. The fraction threshold $f_{\mathrm{thr}}$
is the principal tunable parameter.

\subsection{Per-channel weighting}\label{ssec:weighting}

Vanilla consensus weights every channel equally. We learn per-channel
log-likelihood-ratio weights from labelled training shots:
\[
w_c \;=\; \log \frac{P\bigl(|z_c| > z_{\mathrm{thr}} \mid \text{precursor}\bigr) + \alpha}
                    {P\bigl(|z_c| > z_{\mathrm{thr}} \mid \text{quiescent}\bigr) + \alpha},
\]
with Laplace smoothing $\alpha = 1$ over the binary-event counts. We
clamp $w_c \ge 0$: a channel that fires more in quiescence than in
precursor adds noise rather than signal. The weighted consensus then
replaces $f(t)$:
\[
f_w(t) \;=\; \frac{\sum_c w_c \, \mathbb{1}\bigl[|z_c(t)| > z_{\mathrm{thr}}\bigr]}{\sum_c w_c}.
\]
The weights are learned on training-fold shots only; the test fold
sees only the inferred weights.

\subsection{Multi-scale features}\label{ssec:multiscale}

A single baseline window cannot capture both fast (locked-mode MHD) and
slow (radiation collapse, density limit) precursors. We compute the
per-channel z-score at three scales $B \in \{50, 200, 500\}$ ms, then
derive the following per-time-step features:
\[
\bigl\{f_w(t), R(t), R_{\mathrm{excess}}(t), n_{\mathrm{novel}}(t)\bigr\} \times 3 \text{ scales},
\]
plus two time-derivatives ($\dot{f_w}$ and $\dot{R}_{\mathrm{excess}}$
at the 200 ms scale), for 14 features. Here $R(t)$ is the Kuramoto
order parameter computed from the analytic phases of the channel
z-scores (Hilbert transform), and $R_{\mathrm{excess}}$ is $R(t)$ minus
its trailing rolling median (window 300--600 samples depending on the
scale).

\subsection{Per-channel peak features}\label{ssec:peakz}

We add ten further features: $\max |z_c|$ within a trailing 500 ms
window per channel. This gives the classifier discrimination by
\emph{which} channels are most active, not just how many channels co-fire.

\subsection{Pulse-phase indicators}\label{ssec:phase}

Three binary indicators encode the pulse phase from the
ramp-up/flat-top/ramp-down boundaries in the dataset metadata. These let
the classifier learn phase-conditional thresholds (e.g.\ ``ignore
Kuramoto rise during controlled ramp-down''). Crucially, the training
sampler must include non-disruptive ramp-down windows as negatives to
avoid the phase indicator becoming a label proxy; we document this
explicit fix in \S\ref{ssec:leakage}.

\subsection{Classifier}\label{ssec:classifier}

We use the scikit-learn \texttt{GradientBoostingClassifier} with 250
trees, max-depth 5, learning rate 0.05, subsample 0.7, minimum samples
per leaf 10. Class imbalance ($\sim$71/429 in our windowed sample after
dense extraction) is handled by per-sample weights inversely proportional
to class frequency. The classifier's per-window posterior $P(t)$ feeds
into a final alert rule: $P(t) \ge p_{\mathrm{thr}}$ sustained for
$\tau_{\mathrm{min}}$ samples.

\subsection{Alert masking}\label{ssec:masking}

To avoid alerts firing on plasma ramp-up artefacts or on legitimate
controlled ramp-downs, we restrict the alert evaluation window to
\[
\bigl[t_{\mathrm{flat\_top\_min}} + 50\text{ ms},\; \min(t_{\mathrm{disrupt}},\; t_{\mathrm{flat\_top\_max}} + 20\text{ ms})\bigr].
\]
This baseline-shifted lower bound ensures the rolling baseline window
has filled with flat-top data before any novelty assessment.

\section{Data}\label{sec:data}

We use the MAST Disruption Detection Dataset (Sharma 2025, Zenodo
\texttt{10.5281/zenodo.16032053}). The dataset comprises 500 NetCDF
shot files, each containing 2000 samples at 4638 Hz spanning $\sim$0.43
s of tokamak operation. Ten plasma diagnostics are included per shot:
plasma current ($I_p$), total radiated power, core electron density and
temperature, internal inductance, dynamic loop voltage, two
$D_\alpha$ channels (HM10, HU10), saddle coil channel 6, and
soft-X-ray channel 6. Disruption time, ramp-up boundary, and flat-top
boundaries are provided as NetCDF attributes for the 429 disruptive
shots; the remaining 71 shots are labelled non-disruptive.

The dataset is not balanced: 429 disruptive vs 71 non-disruptive,
roughly 6:1. We mitigate this with class-weight balancing during
classifier training, but the relatively small non-disruptive sample is
the principal limitation on achievable false-alarm-rate estimation.

\section{Iteration log}\label{sec:iterations}

\begin{table}[h]
\centering
\caption{Pareto trajectory across ten iterations. All numbers from 5-fold
CV with the same shot indices. TTD = median time-to-disruption.
Iteration 9 (\textbf{GBM v4 phase}, leak-fixed) is the canonical Pareto
champion. Iteration 10 (lean v5) is a documented negative result.}
\label{tab:trajectory}
\begin{tabular}{c l rrr l}
\toprule
\# & Configuration & TPR (\%) & FAR (\%) & TTD (ms) & Note \\
\midrule
1 & Unweighted consensus (frac=0.50) & 50 & 25 & 19 & Vanilla baseline \\
2 & Weighted (frac=0.30) & 65 & 77 & 65 & High-recall corner \\
2 & Weighted (frac=0.50) & 50 & 25 & 12 & Equivalent to (1) \\
3 & Hybrid K=2 of 3 (A,B,C gates) & 45 & 32 & 16 & \textbf{Negative result} \\
3 & Hybrid K=3 of 3 & 11 & 3 & 6 & High-precision niche \\
4 & Learned LR (7 single-scale feats) & 49 & 14 & 40 & Real gain over hybrid \\
5 & Multi-scale LR (14 feats, 3 scales) & 56--71 & 18--32 & 53--76 & Big lift \\
6 & GBM v1 (14 feats, 150 trees) & 66--81 & 18--44 & 71--102 & GBM dominates LR \\
7 & GBM v2 (dense sampling, 250 trees) & 72--83 & 18--46 & 71--100 & +6 pp TPR at fixed FAR \\
8 & GBM v3 (+per-channel peakz, 24 feats) & 72--82 & 21--45 & 71--100 & Plateau \\
9 & \textbf{GBM v4 (+phase, 27 feats, leak-fixed)} & \textbf{49--73} & \textbf{6--35} & \textbf{65--88} & \textbf{Pareto champion} \\
10 & Lean v5 (9 feats from failure analysis) & 48--75 & 14--39 & 25--39 & \textbf{Negative result} \\
\bottomrule
\end{tabular}
\end{table}

\subsection{Negative result: hybrid pendulum/Kuramoto voting}\label{ssec:hybrid}

Inspired by an internal research note proposing that wave-interference
networks obey nonlinear-pendulum dynamics with golden-ratio equilibria
\cite{knopp_golden_pendulum}, we attempted to combine three structurally
distinct gates into a K-of-3 voting rule:

\begin{itemize}[leftmargin=*,itemsep=0pt]
\item \textbf{Gate A} (coherence-of-derivative): each channel must have
  $|z_c| > 3$ \emph{and} $\mathrm{sign}(\dot{z}_c)$ agreeing with the
  majority of other above-threshold channels.
\item \textbf{Gate B} (Kuramoto): the order parameter $R(t)$ exceeds its
  rolling baseline by a configurable $\Delta R$.
\item \textbf{Gate C} (Hamiltonian rate): $H(t) = \sum_c w_c\, z_c^2(t)$
  has a robust-z-scored $|\dot{H}|$ exceeding threshold.
\end{itemize}

The K=2 vote rule yielded TPR 45\% at FAR 32\%, \emph{worse} than the
weighted single-gate baseline (45/20) at the same TPR. K=3 carved out a
high-precision niche (11/3, precision 96\%) but missed too many
disruptions to serve as a primary detector. Section
\ref{ssec:learned} explains why: on real plasma data the three gates
turned out to be correlated, not independent --- they all detect
``coordinated multi-channel motion'' with different math, and
non-disruptive events with coordinated motion (ELMs, mode-locking,
controlled ramp-downs) fire all three gates simultaneously.

\subsection{Recovery via learned-feature classifier}\label{ssec:learned}

The hybrid failure motivated using the gate scalars as
\emph{input features} to a learned classifier rather than voting on them
directly. The single-scale 7-feature logistic regression (iteration 4)
recovered the loss and added a measurable gain over the weighted
baseline: 49\% TPR at 14\% FAR with 40 ms TTD, a Pareto improvement on
all three metrics simultaneously.

The learned coefficients revealed why the K-of-3 vote failed: gate A's
\texttt{coherence} coefficient was 0.04, gate C's \texttt{Hamiltonian}
coefficient was 0.00, and \texttt{max\_channel\_z} was 0.00. Only
\texttt{R\_excess} (+1.10), \texttt{R} (−1.25, negative because absolute
R is high in both classes), \texttt{weighted\_consensus} (+1.05) and
\texttt{n\_channels\_novel} (−0.57) carried weight. The hybrid had
attempted to vote on signals that contained no orthogonal information.

\subsection{Multi-scale lift}\label{ssec:multilift}

Adding the same features at three baseline scales (50, 200, 500 ms) +
two time-derivatives lifted the LR Pareto: TPR 71\% at FAR 32\% with 76
ms TTD at $p_{\mathrm{thr}} = 0.85$. The 200 ms scale of
\texttt{weighted\_consensus} dominated (LR coefficient +2.03 in the
3-fold average), with 500 ms scales of $R$ and $R_{\mathrm{excess}}$
contributing the next-largest coefficients. Time-derivatives carried
essentially zero weight ($+0.05$ and $+0.01$).

\subsection{Gradient boosting takes over}\label{ssec:gbm}

A gradient-boosting classifier on the same 14 multi-scale features
(iteration 6) further lifted the Pareto: TPR 81\% at FAR 44\% with 102
ms TTD at $p_{\mathrm{thr}} = 0.85$. Dense training-window sampling
(2 ms stride, 50 ms precursor lead, iteration 7) added another $\sim$2 pp
TPR.

\subsection{Phase indicators and a leakage fix}\label{ssec:leakage}

Iteration 9 added three pulse-phase one-hot indicators
(\texttt{in\_rampup}, \texttt{in\_flattop}, \texttt{in\_rampdown}). On
the first attempt, the model achieved TPR 95\% at FAR 92\% --- a
classic feature-leakage signature. The model had discovered that
non-disruptive shots are not sampled past their flat-top end in the
training sampler, so \texttt{in\_rampdown} effectively became a perfect
label proxy for disruptive vs non-disruptive in the windowed sense.

The fix was to extend non-disruptive shot sampling beyond the flat-top
end through the full data range. With the leakage closed, GBM v4
produced the canonical Pareto champion: TPR 73\% at FAR 17\% with 78 ms
TTD at $p_{\mathrm{thr}} = 0.85$ on the seed-42 fold split.

\section{Stability cross-validation}\label{sec:stability}

We re-ran the GBM v4 5-fold CV with three different random seeds for
the fold split (seeds 42, 7, 2024) to confirm the result is robust to
which shots happen to land in which fold. Table \ref{tab:stability}
summarises the per-seed and aggregate statistics.

\begin{table}[h]
\centering
\caption{GBM v4 stability across three CV-fold-split seeds. All values
in per-cent (TPR, FAR) or milliseconds (TTD). The \textbf{p=0.85}
operating point is the canonical Pareto champion; standard deviations
are computed across the three seeds.}
\label{tab:stability}
\begin{tabular}{c c c c}
\toprule
$p_{\mathrm{thr}}$ & TPR mean $\pm$ std & FAR mean $\pm$ std & TTD mean $\pm$ std (ms) \\
\midrule
0.85 & \textbf{74.3 $\pm$ 1.0} & \textbf{14.1 $\pm$ 2.3} & \textbf{80.8 $\pm$ 2.4} \\
0.90 & 66.0 $\pm$ 0.8 & 7.5 $\pm$ 3.7 & 71.4 $\pm$ 1.3 \\
0.95 & 49.3 $\pm$ 0.7 & 3.8 $\pm$ 1.3 & 64.2 $\pm$ 0.5 \\
\bottomrule
\end{tabular}
\end{table}

TPR standard deviations are below 1 pp across seeds; TTD standard
deviations are below 3 ms. The Pareto champion at p=0.85 is
\textbf{74.3 $\pm$ 1.0\% TPR / 14.1 $\pm$ 2.3\% FAR / 80.8 $\pm$ 2.4 ms
TTD}. This is not a fold-split lottery winner.

\section{Failure analysis}\label{sec:failures}

We characterise the 24\% of disruptions missed at the canonical
$p_{\mathrm{thr}} = 0.85$ operating point. For each shot we compute peak
feature values in the analysis window and compare disruptive-shot and
non-disruptive-shot distributions.

\subsection{Discriminating features}

Table \ref{tab:dist} shows median and quartile statistics for the
principal features. The result is sharp:

\begin{table}[h]
\centering
\caption{Per-shot peak feature distributions. Disruptive-class median
versus non-disruptive-class median. Strong discrimination means a
$>2\times$ ratio in the same direction. \emph{Wrong-direction} features
fire \emph{after} disruption, so the windowed pre-disruption pre-image
shows the reverse pattern.}
\label{tab:dist}
\begin{tabular}{l rr l}
\toprule
Feature & disr.\ median & non-disr.\ median & Discrimination \\
\midrule
\texttt{max\_weighted\_consensus\_200ms} & 0.677 & 0.266 & strong (2.5$\times$) \\
\texttt{max\_n\_novel\_500ms} (p75) & 4.0 & 1.0 & strong upper-tail \\
\texttt{max\_R\_500ms} & 0.800 & 0.800 & \emph{none} \\
\texttt{max\_R\_excess\_500ms} & 0.800 & 0.800 & \emph{none} \\
\texttt{saddle\_coil\_peakz} & 0.000 & 0.180 & \emph{wrong direction} \\
\texttt{soft\_xray\_peakz} & 0.000 & 0.598 & \emph{wrong direction} \\
\texttt{ip\_peakz} & 0.000 & 0.569 & \emph{wrong direction} \\
\bottomrule
\end{tabular}
\end{table}

The wrong-direction features are the most physically meaningful
diagnosis. The plasma current $I_p$, the saddle-coil locked-mode
amplitude, and the soft-X-ray emission ARE the canonical disruption
indicators in the plasma-physics literature. However, they deflect
\emph{at} the disruption time, not before. The pre-disruption analysis
window ends at $t_{\mathrm{disrupt}}$, so disruptive-shot windows are
\emph{quieter} on these channels than non-disruptive flat-top windows
(which include normal-physics fluctuations). The classifier therefore
uses these channels in reverse: ``$I_p$ quiet $\Rightarrow$ disruption
imminent.'' This is a fragile signal direction and a key reason for the
remaining gap to RNN baselines, which presumably learn a richer
within-channel shape feature rather than a simple peak-z reduction.

Equally striking: \emph{the Kuramoto order parameter $R$ itself does
not discriminate.} Both classes routinely reach $R \approx 0.8$. Only
the \emph{rise of $R$ over a rolling baseline} (i.e.,\ $R_{\mathrm{excess}}$)
carries information.

\subsection{Structurally invisible disruptions}

Defining a disruption as ``structurally invisible'' if all its peak
feature values lie below the 90th percentile of non-disruptive shots,
we find $8 / 417 = 2\%$ of disruptions in this regime. The remaining
miss budget at $p_{\mathrm{thr}} = 0.85$ is therefore a
threshold-tuning problem, not a missing-information problem.

\subsection{Negative experiment: the lean feature set}\label{ssec:lean}

The failure analysis motivated a 9-feature lean redesign: drop the
no-discrimination $R$-absolute features and the wrong-direction peakz
features, add \texttt{weighted\_consensus} dynamics (rate of rise,
duration above 0.3, integral over 500 ms).

Result: TPR 48\% at FAR 14\% with 25 ms TTD at $p_{\mathrm{thr}} = 0.90$
--- much worse than GBM v4 (74/14/81 at $p_{\mathrm{thr}} = 0.85$). The
``noisy'' features the failure analysis had identified as candidates
for removal apparently carried \emph{contextual} information the GBM
used effectively. Stripping them sacrificed the model's ability to
disambiguate in specific pulse-phase contexts.

The lesson is that GBM feature-importance rankings are not a clean
guide to which features can be discarded: a feature with low marginal
importance may still be load-bearing within a specific tree split.
Iteration 10 stands as documented evidence of this.

\section{Comparison with published baselines}\label{sec:baselines}

The Ferreira et al.\ 2019 \cite{ferreira2019} LSTM on JET bolometer
data reports approximately 85\% TPR at 15\% FAR. Our final operating
point is 74.3\% TPR at 14.1\% FAR --- approximately matching the
RNN false-alarm-rate while running approximately 11 pp below the
RNN's TPR.

Three caveats apply to this comparison:

\begin{enumerate}[leftmargin=*,itemsep=0pt]
\item Different machine (MAST versus JET). MAST disruptions are dominated
  by density-limit and radiation-collapse precursors; JET disruptions
  also include locked-mode and tearing-mode causes. Per-machine
  precursor distributions differ.
\item Different dataset size (500 versus $\sim$1500 shots in Ferreira's
  evaluation set).
\item Different baseline pipelines. Ferreira uses raw bolometer profiles
  fed to an LSTM, learning temporal shape implicitly; we use hand-crafted
  scalar features and gradient boosting, with all temporal structure
  collapsed to peak/rate/integral summaries.
\end{enumerate}

Closing the remaining 11 pp TPR gap most likely requires one or more of:

\begin{itemize}[leftmargin=*,itemsep=0pt]
\item \textbf{Operational metadata for multi-class precursor learning.}
  Disruption-cause labels (radiation, density limit, locked mode,
  H/L back-transition, $\ldots$) would let a multi-class model learn
  cause-specific precursor signatures rather than a single
  binary-disruption signature. The current MAST dataset does not
  expose these labels.
\item \textbf{Larger training set.} 429 disruptive shots is small for
  the diversity of physics involved.
\item \textbf{Sequence-aware architectures.} An LSTM or transformer
  could learn within-channel \emph{shape} features that our per-window
  scalar features collapse out. The wrong-direction peakz features
  identified in \S\ref{sec:failures} would be unaffected because the
  network would access the time series before disruption directly.
\end{itemize}

\section{Reproducibility}\label{sec:repro}

All code, dataset URLs, per-iteration results JSON, and commit history
are available at \url{https://github.com/Zynerji/JET-Disruption}. The
canonical Pareto-champion result is reproducible from a clean Python
3.12 environment with
\texttt{pip install numpy scipy scikit-learn xarray netCDF4} and:

\begin{verbatim}
git clone https://github.com/Zynerji/JET-Disruption
cd JET-Disruption
python -m pip install netCDF4 xarray scikit-learn
python -c "import urllib.request; urllib.request.urlretrieve(
    'https://zenodo.org/api/records/16032053/files/T_lead_0ms.tar.xz/content',
    'data/T_lead_0ms.tar.xz')"
mkdir -p data && tar -xJf data/T_lead_0ms.tar.xz -C data
python examples/benchmark_mast_gbm_v4.py
python examples/stability_check_v4.py
\end{verbatim}

Each benchmark writes a results JSON next to itself. The full Pareto
trajectory in Table~\ref{tab:trajectory} is the literal output of
running iterations 1 through 10 in sequence.

\section{Conclusion}\label{sec:conclusion}

We have demonstrated that the Systroph\=e address-space cross-observable
consensus catcher, originally developed for gravitational-wave search
and for spacetime-metric special-locus identification, applies to
tokamak disruption prediction on real MAST plasma data with measurable
performance. The final operating point of 74.3 $\pm$ 1.0\% TPR at 14.1
$\pm$ 2.3\% FAR with 80.8 $\pm$ 2.4 ms median time-to-disruption is
within 3 pp of the published RNN false-alarm rate, $\sim$11 pp below
the RNN true-positive rate, and is statistically stable across
fold-split seeds.

The work is documented as an ``honest iteration log'' with two outright
negative results (hybrid voting; failure-analysis-informed feature
reduction). The negative results were as informative as the positive
ones: they constrained subsequent design choices and exposed a
systematic feature-leakage failure mode in the pulse-phase indicator
that would have invalidated naive bench-marking.

The remaining gap to RNN baselines is, we argue, in operational-metadata
access and sequence-aware shape learning rather than in the consensus
methodology itself. The methodology generalises cleanly across
domains --- gravitational waves, spacetime physics, plasma disruptions
--- which we take as evidence that cross-observable consensus is a
fundamentally good detector design for regime-change problems.

\begin{thebibliography}{9}

\bibitem{vega2013} Vega, J., et al.\ \emph{Results of the JET real-time disruption
predictor in the ITER-like wall campaigns.} Fusion Eng.\ Des.\ \textbf{88},
1228--1231 (2013).

\bibitem{maraschek2018} Maraschek, M., et al.\ \emph{Path-oriented early
reaction to approaching disruptions in ASDEX Upgrade and TCV.} Plasma
Phys.\ Control.\ Fusion \textbf{60}, 014047 (2018).

\bibitem{ferreira2019} Ferreira, D.R., Carvalho, P.J., Fernandes, H.
\emph{Deep Learning for Plasma Tomography and Disruption Prediction
from Bolometer Data.} IEEE Trans.\ Plasma Sci.\ \textbf{48}, 2--12 (2020).

\bibitem{sharma2025} Sharma, P.\ \emph{MAST disruption detection
dataset.} Zenodo, doi:10.5281/zenodo.16032053 (2025).

\bibitem{systrophe_gw} Knopp, C.\ \emph{Systroph\=e v2 GW unmodeled-burst
catcher: per-detector excess-sigma + coherence + band-ratio with
2-of-3 coincidence.} Systroph\=e v0.19.2 (2026). See \texttt{src/systrophe/catcher\_v2.py}
and \texttt{src/systrophe/gw\_pipeline\_v3.py} at
\url{https://github.com/Zynerji/systrophe}.

\bibitem{systrophe_lp_ridge} Knopp, C.\ \emph{Cross-observable consensus
localizes the Lewis--Papapetrou chronology horizon.} Systroph\=e v0.19.2
whitepaper, \texttt{paper/lp\_ridge\_localization.pdf} (2026).

\bibitem{knopp_golden_pendulum} Knopp, C.\ \emph{The Golden Pendulum:
Deriving Conservation Laws in Wave Interference Networks.} Internal
draft (Jan 2026), referenced here only for the methodological hybrid-gate
inspiration; the conservation-law claims of that draft are not
validated by the present work.

\end{thebibliography}

\end{document}
